% DONE 
%   Provide a concise description and comments on your algorithmic design and implementation
%   Explain the main implemented functionality, key parameters and conceptual aspects of your     program.
%   Example: explain how you defined the cost term and LiDAR scan processing for DWA, etc.
%   Briefly comment the main structure of your ROS 2 program (do not copy your code in the report,  just highlights the main elements and the workflow: nodes, publishers/subscribers to relevant topics and parameters)

\section{Design}
\label{sec:design}

\subsection{System Architecture}
The navigation system is implemented as a ROS 2 node (\texttt{controller\_node}) that integrates the Dynamic Window Approach (DWA) algorithm for obstacle avoidance and dynamic goal tracking. The system operates at a fixed control frequency of 15 Hz, managed by a ROS 2 timer callback that executes the main control loop.

The system is designed to work both in simulation (using Gazebo with TurtleBot3) and on the real robot platform, with conditional logic handling the differences in goal detection methods: odometry-based goals in simulation versus AprilTag-based visual detection on the physical robot.

\subsection{ROS 2 Node Structure}

\subsubsection{Publishers and Subscribers}
The controller node establishes communication through the following main topics:

\textbf{Subscribers:}
\begin{itemize}
    \item \texttt{/odom} (Odometry): Receives robot pose and velocity estimates
    \item \texttt{/scan} (LaserScan): Receives LiDAR measurements for obstacle detection
    \item \texttt{/dynamic\_goal\_pose} (Odometry): Receives goal positions in simulation mode
    \item \texttt{/camera/landmarks} (LandmarkArray): Receives AprilTag detections on real robot
\end{itemize}

\textbf{Publishers:}
\begin{itemize}
    \item \texttt{/cmd\_vel} (Twist): Sends velocity commands to the robot
    \item \texttt{/goal\_marker} (Marker): Used to visualize the current goal in RViz
    \item \texttt{/dwa\_feedback} (String): Provides navigation task feedback
    \item \texttt{/goal\_reached}, \texttt{/collision\_flag} (Bool): Publishes task status flags
    \item \texttt{/filter\_scan} (MarkerArray): Visualizes filtered LiDAR points
    \item \texttt{/task\_status} (String): Publishes overall task status (e.g. "Goal", "Collision", "Timeout")
    \item \texttt{/goal\_pose} (Point): Republishes the current goal position for visalization purposes (in odom frame)
\end{itemize}

\subsubsection{State Management}
The node maintains two distinct state representations:
\begin{itemize}
    \item \textbf{Global state} $[x_{global}, y_{global}, \theta_{global}, v, \omega]$: Updated from odometry, used for coordinate transformations
    \item \textbf{Robot-centric state} $[0, 0, 0, v, \omega]$: Always centered at the robot's origin, used directly by DWA for trajectory planning
\end{itemize}

This dual-state design simplifies the DWA implementation by maintaining goals and obstacles in the robot's reference frame, eliminating the need for continuous coordinate transformations during trajectory evaluation.

\subsection{LiDAR Scan Processing}

The LiDAR data processing pipeline implements several filtering stages to ensure robust obstacle detection while maintaining computational efficiency:

\begin{enumerate}
    \item \textbf{Range validation}: Invalid measurements (NaN, Inf) are handled by assigning maximum range (3.5 m) to avoid false obstacle detections
    \item \textbf{Distance filtering}: Only measurements within the effective range [0.12 m, 3.5 m] are considered for obstacle detection
    \item \textbf{Angular sectoring}: The full 360° scan (could be less in real robot) is divided into $N_{ranges}$ angular sectors to reduce computational load
    \item \textbf{Minimum distance selection}: Within each sector, only the closest obstacle point is retained, representing the most critical obstacle in that direction
    \item \textbf{Coordinate transformation}: Range-bearing measurements are converted to Cartesian coordinates $(x, y)$ in the robot's reference frame
\end{enumerate}

This filtering approach reduces the number of obstacle points to $N_{ranges}$, significantly improving computational performance while maintaining sufficient environmental representation for safe navigation.

\subsection{Goal Management}

The goal manager implements dynamic target switching capabilities, allowing the robot to update its navigation objective at any time during task execution. Two distinct modes are implemented:

\textbf{Simulation mode:} Goals are received as odometry messages and transformed from the global frame to the robot's local frame through rotation by $-\theta$:
\begin{equation}
    \begin{bmatrix} x_{robot} \\ y_{robot} \end{bmatrix} = 
    \begin{bmatrix} \cos(-\theta) & -\sin(-\theta) \\ \sin(-\theta) & \cos(-\theta) \end{bmatrix}
    \begin{bmatrix} x_{goal} - x_{robot,global} \\ y_{goal} - y_{robot,global} \end{bmatrix}
\end{equation}

\textbf{Real robot mode:} AprilTag detections provide range-bearing measurements that are directly converted to robot-frame coordinates:
\begin{equation}
    \begin{bmatrix} x_{goal} \\ y_{goal} \end{bmatrix} = 
    \begin{bmatrix} range \cdot \cos(bearing) \\ range \cdot \sin(bearing) \end{bmatrix}
\end{equation}

The goal is deleted (\texttt{None}) if no landmark is detected, preventing the robot from pursuing outdated targets.

\subsection{Safety Mechanisms}

Multiple safety layers are implemented to ensure collision-free navigation:

\begin{enumerate}
    \item \textbf{Collision detection}: If any LiDAR measurement falls below the collision tolerance, the robot immediately stops and sets a collision flag. If the obstacle is no longer detected, navigation resumes.
    \item \textbf{DWA trajectory scoring}: Trajectories that would bring the robot closer than $r_{robot} + d_{collision}$ (where $r_{robot} = 0.2$ m) to any obstacle receive a heavy penalty ($-100$) in the objective function
    \item \textbf{Timeout protection}: Navigation tasks exceeding 300 control steps are automatically terminated
    \item \textbf{Sensor validation}: Commands are only computed when both LiDAR data and a valid goal are available
\end{enumerate}

\subsection{DWA Implementation and Objective Functions}

The DWA controller generates trajectories by sampling the dynamic window of feasible velocities and evaluating them according to an objective function. Three variants are implemented:

\subsubsection{Task 1: Original DWA}
\label{sssec:dwa_cost}
The standard objective function (hereafter referred to as \texttt{obj1}) balances three terms:
\begin{equation}
    J = \alpha \cdot heading + \beta \cdot vel + \gamma \cdot dist_{obst}
\end{equation}

where:
\begin{itemize}
    \item \textbf{heading}: measures alignment between trajectory endpoint and goal direction.
    \item \textbf{vel}: promotes forward motion by scoring linear velocity, enhanced by an exponential term that increases as the robot approaches the goal
    \item \textbf{dist\_obst}: rewards maintaining distance from obstacles, with collision penalties for unsafe trajectories
\end{itemize}

\subsubsection{Task 2a: Velocity Reduction Near Goal}
\label{sssec:task2a}
This variant (hereafter referred to as \texttt{obj2a}) introduces an additional penalty term to slow down the robot when approaching the goal:
\begin{equation}
    vel' = -v \cdot \left(1 - \frac{d_{goal}}{d_{threshold}}\right), \quad \text{if } d_{goal} < d_{threshold}
\end{equation}

where $d_{threshold}$ m defines the activation range. The penalty increases linearly as the robot gets closer to the goal, effectively reducing the selected velocities in the final approach phase.

\subsubsection{Task 2b: Target Following}
\label{sssec:task2b}
For dynamic target following (hereafter referred to as \texttt{obj2b}), a distance maintenance term is added:
\begin{equation}
    J = \alpha \cdot heading + \beta \cdot vel + \gamma \cdot dist_{obst} + \delta \cdot dist_{target}
\end{equation}

The $dist_{target}$ term is computed as:
\begin{equation}
    dist_{target} = \exp\left(-\frac{|d_{actual} - d_{desired}|}{d_{desired}}\right)
\end{equation}

Additional penalties (50\% reduction if $d < 0.5 \cdot d_{desired}$, 70\% if $d > 2.0 \cdot d_{desired}$) discourage getting too close or too far from the target, ensuring optimal visual tracking conditions.

\subsection{Control Loop Execution}

The main control callback implements the following execution flow:

\begin{enumerate}
    \item Check goal reach condition: $\|goal_{robot}\| < d_{tol}$
    \item Validate sensor data availability and goal validity
    \item Check for collision, timeout, or goal reached conditions
    \item If all checks pass: compute velocity command using DWA
    \item Publish velocity command, status flags, and periodic feedback
    \item Increment control step counter
\end{enumerate}

When any termination condition is met (goal reached, collision, timeout, or no goal), the robot publishes a zero velocity command to ensure safe stopping.

